\input{preamble}

% OK, start here.
%
\begin{document}

\title{Topologies sur les espaces algébriques}

\maketitle

\phantomsection
\label{section-phantom}

\tableofcontents




\section{Introduction}
\label{section-introduction}

\noindent
Dans ce chapitre, nous introduisons quelques topologies sur la
catégorie des espaces algébriques. Comparer avec le contenu de \cite{SGA1},
\cite{Ner}, \cite{LM-B} et \cite{Kn}.
Avant cela, signalons que
de nombreux choix de sites (au sens de
Sites, définition \ref{sites-definition-site}) donnent
la même notion de faisceau sur la catégorie sous-jacente. Ainsi,
nos choix peuvent différer légèrement de ceux des références,
mais conduisent finalement aux mêmes groupes de cohomologie, etc.




\section{Le procédé général}
\label{section-procedure}

\noindent
Dans cette section, nous expliquons un procédé général pour construire les
sites avec lesquels nous travaillerons. Cette discussion n'aura guère de
sens si le lecteur n'a pas lu
Topologies, section \ref{topologies-section-procedure}.

\medskip\noindent
Soit $S$ un schéma de base.
Prenons une catégorie $\Sch_\alpha$ construite comme dans
Ensembles, lemme \ref{sets-lemma-construct-category} en partant de
$S$ et de tout ensemble de schémas sur $S$ que l'on souhaite inclure.
Choisissons un ensemble quelconque de
recouvrements $\text{Cov}_{fppf}$ sur $\Sch_\alpha$, comme dans
Ensembles, lemme \ref{sets-lemma-coverings-site},
en partant de la catégorie $\Sch_\alpha$ et de la classe des recouvrements fppf.
Notons $\Sch_{fppf}$ le grand site fppf ainsi
obtenu, et notons $(\Sch/S)_{fppf}$ le
grand site fppf de $S$. (Tout ceci est exactement prescrit dans Topologies,
section \ref{topologies-section-fppf}).

\medskip\noindent
Pour les choix ci-dessus, la catégorie des espaces algébriques sur $S$
possède un ensemble de classes d'isomorphisme. Pour le voir, on peut utiliser le
fait que tout espace algébrique sur $S$ est de la forme $U/R$ pour
une certaine relation d'équivalence \'etale $j : R \to U \times_S U$ avec
$U, R \in \Ob((\Sch/S)_{fppf})$, voir
Espaces, lemme \ref{spaces-lemma-space-presentation}.
On peut donc trouver une sous-catégorie pleine $\textit{Spaces}/S$ de la catégorie des
espaces algébriques sur $S$ qui possède un ensemble d'objets
tel que tout espace algébrique soit isomorphe à un objet de
$\textit{Spaces}/S$. Nous fixons une telle catégorie.

\medskip\noindent
Dans les sections suivantes, pour une topologie $\tau$ donnée, le grand site
$(\textit{Spaces}/S)_\tau$ (resp.\ le grand site $(\textit{Spaces}/X)_\tau$
d'un espace algébrique $X$ sur $S$)
a pour catégorie sous-jacente la catégorie $\textit{Spaces}/S$
(resp.\ la sous-catégorie $\textit{Spaces}/X$ de $\textit{Spaces}/S$, voir
Catégories, exemple \ref{categories-example-category-over-X}).
Le procédé pour en faire un site consiste, comme d'habitude, à définir une
classe de recouvrements $\tau$ et à utiliser
Ensembles, lemme \ref{sets-lemma-coverings-site},
pour choisir un ensemble suffisamment grand de recouvrements définissant la topologie.

\medskip\noindent
Signalons que le {\it petit site \'etale $X_\etale$
d'un espace algébrique $X$} a déjà été défini dans
Propriétés des espaces, définition
\ref{spaces-properties-definition-etale-site}.
Ses objets sont les schémas \'etales sur $X$, qui sont nombreux
par définition d'un espace algébrique. Cependant,
un site plus naturel du point de vue de ce chapitre (comparer
Topologies, définition \ref{topologies-definition-big-small-etale})
est le site $X_{spaces, \etale}$ de
Propriétés des espaces, définition
\ref{spaces-properties-definition-spaces-etale-site}.
Ces deux sites définissent le même topos, voir
Propriétés des espaces, lemme \ref{spaces-properties-lemma-compare-etale-sites}.
Nous ne les redéfinirons pas dans ce chapitre ; nous les
utiliserons simplement.






\section{Topologie de Zariski}
\label{section-zariski}

\noindent
Dans
Espaces, section \ref{spaces-section-Zariski}
nous avons introduit la notion de recouvrement de Zariski d'un espace algébrique par
des sous-espaces ouverts. Voici la notion correspondante, où les sous-espaces ouverts
sont remplacés par des immersions ouvertes.

\begin{definition}
\label{definition-zariski-covering}
Soit $S$ un schéma et soit $X$ un espace algébrique sur $S$.
Un {\it recouvrement de Zariski de $X$} est une famille de morphismes
$\{f_i : X_i \to X\}_{i \in I}$ d'espaces algébriques sur $S$
telle que chaque $f_i$ soit une immersion ouverte
et telle que
$$
|X| = \bigcup\nolimits_{i \in I} |f_i|(|X_i|),
$$
c'est-à-dire que la famille des morphismes est surjective.
\end{definition}

\noindent
Bien que les recouvrements de Zariski soient parfois utiles, la topologie correspondante
sur la catégorie des espaces algébriques est beaucoup trop grossière et peu
utile. Elle définit néanmoins un site.

\begin{lemma}
\label{lemma-zariski}
Soit $S$ un schéma.
Soit $X$ un espace algébrique sur $S$.
\begin{enumerate}
\item Si $X' \to X$ est un isomorphisme, alors $\{X' \to X\}$
est un recouvrement de Zariski de $X$.
\item Si $\{X_i \to X\}_{i\in I}$ est un recouvrement de Zariski et si, pour chaque
$i$, nous avons un recouvrement de Zariski $\{X_{ij} \to X_i\}_{j\in J_i}$, alors
$\{X_{ij} \to X\}_{i \in I, j\in J_i}$ est un recouvrement de Zariski.
\item Si $\{X_i \to X\}_{i\in I}$ est un recouvrement de Zariski
et si $X' \to X$ est un morphisme d'espaces algébriques, alors
$\{X' \times_X X_i \to X'\}_{i\in I}$ est un recouvrement de Zariski.
\end{enumerate}
\end{lemma}

\begin{proof}
Démonstration omise.
\end{proof}










\section{Topologie \'etale}
\label{section-etale}

\noindent
Dans cette section, nous étudions la notion de recouvrement \'etale des
espaces algébriques et définissons le grand site \'etale d'un
espace algébrique. Comparer avec
Topologies, section \ref{topologies-section-etale}.

\begin{definition}
\label{definition-etale-covering}
Soit $S$ un schéma et soit $X$ un espace algébrique sur $S$.
Un {\it recouvrement \'etale de $X$} est une famille de morphismes
$\{f_i : X_i \to X\}_{i \in I}$ d'espaces algébriques sur $S$
telle que chaque $f_i$ soit \'etale
et telle que
$$
|X| = \bigcup\nolimits_{i \in I} |f_i|(|X_i|),
$$
c'est-à-dire que la famille des morphismes est surjective.
\end{definition}

\noindent
C'est exactement la même notion que celle de
Topologies, définition \ref{topologies-definition-etale-covering}.
En particulier, si $X$ et tous les $X_i$ sont des schémas, on retrouve la
notion usuelle de recouvrement \'etale de schémas.

\begin{lemma}
\label{lemma-zariski-etale}
Tout recouvrement de Zariski est un recouvrement \'etale.
\end{lemma}

\begin{proof}
Cela découle immédiatement des définitions et du fait qu'une immersion ouverte
est un morphisme \'etale (cela résulte de
Morphismes, lemme \ref{morphisms-lemma-open-immersion-etale}, via
Espaces, lemme
\ref{spaces-lemma-representable-transformations-property-implication}
puisque les immersions sont représentables).
\end{proof}

\begin{lemma}
\label{lemma-etale}
Soit $S$ un schéma.
Soit $X$ un espace algébrique sur $S$.
\begin{enumerate}
\item Si $X' \to X$ est un isomorphisme, alors $\{X' \to X\}$
est un recouvrement \'etale de $X$.
\item Si $\{X_i \to X\}_{i\in I}$ est un recouvrement \'etale et si, pour chaque
$i$, nous avons un recouvrement \'etale $\{X_{ij} \to X_i\}_{j\in J_i}$, alors
$\{X_{ij} \to X\}_{i \in I, j\in J_i}$ est un recouvrement \'etale.
\item Si $\{X_i \to X\}_{i\in I}$ est un recouvrement \'etale
et si $X' \to X$ est un morphisme d'espaces algébriques, alors
$\{X' \times_X X_i \to X'\}_{i\in I}$ est un recouvrement \'etale.
\end{enumerate}
\end{lemma}

\begin{proof}
Démonstration omise.
\end{proof}

\noindent
Le lemme suivant montre que les sites
$(\textit{Spaces}/X)_\etale$ et $(\textit{Spaces}/X)_{smooth}$
ont la même catégorie de faisceaux.

\begin{lemma}
\label{lemma-etale-dominates-smooth}
Soit $S$ un schéma et soit $X$ un espace algébrique sur $S$.
Soit $\{X_i \to X\}_{i \in I}$ un recouvrement lisse de $X$.
Alors il existe un recouvrement \'etale $\{U_j \to X\}_{j \in J}$
de $X$ qui raffine $\{X_i \to X\}_{i \in I}$.
\end{lemma}

\begin{proof}
Choisissons d'abord un schéma $U$ et un morphisme \'etale surjectif $U \to X$.
Pour chaque $i$, choisissons un schéma $W_i$ et un morphisme \'etale surjectif
$W_i \to X_i$. Alors $\{W_i \to X\}_{i \in I}$ est un recouvrement lisse
qui raffine $\{X_i \to X\}_{i \in I}$. Ainsi,
$\{W_i \times_X U \to U\}_{i \in I}$ est un recouvrement lisse de schémas.
Par le lemme \ref{more-morphisms-lemma-etale-dominates-smooth} de Compléments sur les morphismes,
nous pouvons choisir un recouvrement \'etale $\{U_j \to U\}$ qui raffine
$\{W_i \times_X U \to U\}$. Alors $\{U_j \to X\}_{j \in J}$
est un recouvrement \'etale qui raffine $\{X_i \to X\}_{i \in I}$.
\end{proof}

\begin{definition}
\label{definition-big-etale-site}
Soit $S$ un schéma. Un grand site \'etale {\it $(\textit{Spaces}/S)_\etale$}
est un site construit comme suit :
\begin{enumerate}
\item Choisissons un grand site \'etale $(\Sch/S)_\etale$ comme dans
Topologies, section \ref{topologies-section-etale}.
\item Prenons pour catégorie sous-jacente la catégorie $\textit{Spaces}/S$
des espaces algébriques sur $S$ (voir la discussion de la
section \ref{section-procedure} expliquant qu'il s'agit d'un ensemble).
\item Choisissons un ensemble quelconque de recouvrements comme dans
Ensembles, lemme \ref{sets-lemma-coverings-site}, en partant de la
catégorie $\textit{Spaces}/S$ et de la classe des recouvrements \'etales
de la définition \ref{definition-etale-covering}.
\end{enumerate}
\end{definition}

\noindent
Après cette définition, on peut localiser pour obtenir le site \'etale
d'un espace algébrique.

\begin{definition}
\label{definition-big-small-etale}
Soit $S$ un schéma. Soit $(\textit{Spaces}/S)_\etale$ comme dans
la définition \ref{definition-big-etale-site}.
Soit $X$ un espace algébrique sur $S$, c'est-à-dire un objet de
$(\textit{Spaces}/S)_\etale$. Alors le grand site \'etale
{\it $(\textit{Spaces}/X)_\etale$} de $X$
est la localisation du site $(\textit{Spaces}/S)_\etale$
en $X$, introduite dans Sites, section \ref{sites-section-localize}.
\end{definition}

\noindent
Rappelons que, pour un espace algébrique $X$ sur $S$ comme dans
la définition, nous avons déjà défini les petits sites \'etales
$X_{spaces, \etale}$ et $X_\etale$, voir
Propriétés des espaces, section \ref{spaces-properties-section-etale-site}.
Nous identifierons tacitement les topos correspondants au moyen du
foncteur d'inclusion $X_\etale \subset X_{spaces, \etale}$
(Propriétés des espaces, lemme \ref{spaces-properties-lemma-compare-etale-sites})
et l'appellerons le petit topos \'etale de $X$.
Nous établissons ensuite des relations entre les topos
associés à ces sites.

\begin{lemma}
\label{lemma-put-in-T-etale}
Soit $S$ un schéma. Soit $f : Y \to X$ un morphisme dans
$(\textit{Spaces}/S)_\etale$. Le foncteur d'inclusion
$Y_{spaces, \etale} \to (\textit{Spaces}/X)_\etale$
est cocontinu et induit un morphisme de topos
$$
i_f :
\Sh(Y_\etale)
\longrightarrow
\Sh((\textit{Spaces}/X)_\etale)
$$
Pour un faisceau $\mathcal{G}$ sur $(\textit{Spaces}/X)_\etale$
nous avons la formule $(i_f^{-1}\mathcal{G})(U/Y) = \mathcal{G}(U/X)$.
Le foncteur $i_f^{-1}$ possède également un adjoint à gauche $i_{f, !}$ qui commute
aux produits fibrés et aux égaliseurs.
\end{lemma}

\begin{proof}
Notons $u : Y_{spaces, \etale} \to (\textit{Spaces}/X)_\etale$ le foncteur.
Autrement dit, étant donné un morphisme \'etale $j : U \to Y$ correspondant
à un objet de $Y_{spaces, \etale}$, posons
$u(U \to T) = (f \circ j : U \to S)$.
La catégorie $Y_{spaces, \etale}$
possède des produits fibrés et des égaliseurs, et $u$ commute à ces constructions.
Il est immédiat que $u$ est cocontinu.
Le foncteur $u$ est également continu, car $u$ transforme les recouvrements en recouvrements et
commute aux produits fibrés. Le lemme découle donc de
Sites, lemmes \ref{sites-lemma-when-shriek}
et \ref{sites-lemma-preserve-equalizers}.
\end{proof}

\begin{lemma}
\label{lemma-at-the-bottom-etale}
Soit $S$ un schéma. Soit $X$ un objet de $(\textit{Spaces}/S)_\etale$.
Le foncteur d'inclusion $X_{spaces, \etale} \to (\textit{Spaces}/X)_\etale$
satisfait aux hypothèses de Sites, lemme \ref{sites-lemma-bigger-site},
et induit donc un morphisme de sites
$$
\pi_X :
(\textit{Spaces}/X)_\etale
\longrightarrow
X_{spaces, \etale}
$$
et un morphisme de topos
$$
i_X :
\Sh(X_\etale)
\longrightarrow
\Sh((\textit{Spaces}/X)_\etale)
$$
tel que $\pi_X \circ i_X = \text{id}$. De plus, $i_X = i_{\text{id}_X}$
avec $i_{\text{id}_X}$ comme dans le lemme \ref{lemma-put-in-T-etale}.
En particulier, le foncteur $i_X^{-1} = \pi_{X, *}$ est décrit par la règle
$i_X^{-1}(\mathcal{G})(U/X) = \mathcal{G}(U/X)$.
\end{lemma}

\begin{proof}
Dans ce cas, le foncteur
$u : X_{spaces, \etale} \to (\textit{Spaces}/X)_\etale$,
en plus des propriétés établies dans la démonstration du
lemme \ref{lemma-put-in-T-etale} ci-dessus, est également pleinement fidèle
et transforme l'objet final en objet final.
Le lemme découle de Sites, lemme \ref{sites-lemma-bigger-site}.
\end{proof}

\begin{definition}
\label{definition-restriction-small-etale}
Dans la situation du lemme \ref{lemma-at-the-bottom-etale},
le foncteur $i_X^{-1} = \pi_{X, *}$ est souvent
appelé la {\it restriction au petit site \'etale}, et, pour un faisceau
$\mathcal{F}$ sur le grand site \'etale, nous notons souvent
$\mathcal{F}|_{X_\etale}$ cette restriction.
\end{definition}

\noindent
Avec cette notation, nous avons, pour un faisceau $\mathcal{F}$ sur le
grand site et un faisceau $\mathcal{G}$ sur le petit site,
\begin{align*}
\Mor_{\Sh(X_\etale)}(
\mathcal{F}|_{X_\etale},
\mathcal{G})
& =
\Mor_{\Sh((\textit{Spaces}/X)_\etale)}(
\mathcal{F},
i_{X, *}\mathcal{G}) \\
\Mor_{\Sh(X_\etale)}(
\mathcal{G},
\mathcal{F}|_{X_\etale})
& =
\Mor_{\Sh((\textit{Spaces}/X)_\etale)}(
\pi_X^{-1}\mathcal{G},
\mathcal{F})
\end{align*}
De plus, $(i_{X, *}\mathcal{G})|_{X_\etale} = \mathcal{G}$
et $(\pi_X^{-1}\mathcal{G})|_{X_\etale} = \mathcal{G}$.

\begin{lemma}
\label{lemma-morphism-big-etale}
Soit $S$ un schéma. Soit $f : Y \to X$ un morphisme dans
$(\textit{Spaces}/S)_\etale$. Le foncteur
$$
u :
(\textit{Spaces}/Y)_\etale
\longrightarrow
(\textit{Spaces}/X)_\etale,
\quad
V/Y \longmapsto V/X
$$
est cocontinu et possède un adjoint à droite continu
$$
v :
(\textit{Spaces}/X)_\etale
\longrightarrow
(\textit{Spaces}/Y)_\etale,
\quad
(U \to X) \longmapsto (U \times_X Y \to Y).
$$
Ils induisent le même morphisme de topos
$$
f_{big} :
\Sh((\textit{Spaces}/Y)_\etale)
\longrightarrow
\Sh((\textit{Spaces}/X)_\etale)
$$
Nous avons $f_{big}^{-1}(\mathcal{G})(U/Y) = \mathcal{G}(U/X)$.
Nous avons $f_{big, *}(\mathcal{F})(U/X) = \mathcal{F}(U \times_X Y/Y)$.
De plus, $f_{big}^{-1}$ possède un adjoint à gauche $f_{big!}$ qui commute aux
produits fibrés et aux égaliseurs.
\end{lemma}

\begin{proof}
Le foncteur $u$ est cocontinu, continu et commute aux produits fibrés
et aux égaliseurs (détails omis ; comparer avec la démonstration du
lemme \ref{lemma-put-in-T-etale}).
Par conséquent,
les lemmes de Sites \ref{sites-lemma-when-shriek} et
\ref{sites-lemma-preserve-equalizers}
s'appliquent et nous obtenons la formule
pour $f_{big}^{-1}$ ainsi que l'existence de $f_{big!}$. De plus,
le foncteur $v$ est un adjoint à droite car, pour $U/Y$ et $V/X$,
nous avons $\Mor_X(u(U), V) = \Mor_Y(U, V \times_X Y)$
comme souhaité. Nous pouvons donc appliquer
les lemmes de Sites \ref{sites-lemma-have-functor-other-way} et
\ref{sites-lemma-have-functor-other-way-morphism} pour obtenir la
formule de $f_{big, *}$.
\end{proof}

\begin{lemma}
\label{lemma-morphism-big-small-etale}
Soit $S$ un schéma. Soit $f : Y \to X$ un morphisme dans
$(\textit{Spaces}/S)_\etale$.
\begin{enumerate}
\item Nous avons $i_f = f_{big} \circ i_T$, avec $i_f$ comme dans
le lemme \ref{lemma-put-in-T-etale} et $i_T$ comme dans
le lemme \ref{lemma-at-the-bottom-etale}.
\item Le foncteur $X_{spaces, \etale} \to T_{spaces, \etale}$,
$(U \to X) \mapsto (U \times_X Y \to Y)$ est continu et induit
un morphisme de sites
$$
f_{spaces, \etale} : Y_{spaces, \etale} \longrightarrow X_{spaces, \etale}
$$
Le morphisme correspondant entre les petits topos \'etales est noté
$$
f_{small} : \Sh(Y_\etale) \to \Sh(X_\etale)
$$
Nous avons $f_{small, *}(\mathcal{F})(U/X) = \mathcal{F}(U \times_X Y/Y)$.
\item Nous avons un diagramme commutatif de morphismes de sites
$$
\xymatrix{
Y_{spaces, \etale} \ar[d]_{f_{spaces, \etale}} &
(\textit{Spaces}/Y)_\etale \ar[d]^{f_{big}} \ar[l]^-{\pi_Y}\\
X_{spaces, \etale} &
(\textit{Spaces}/X)_\etale \ar[l]_-{\pi_X}
}
$$
de sorte que $f_{small} \circ \pi_Y = \pi_X \circ f_{big}$ comme morphismes de topos.
\item Nous avons $f_{small} = \pi_X \circ f_{big} \circ i_Y = \pi_X \circ i_f$.
\end{enumerate}
\end{lemma}

\begin{proof}
L'égalité $i_f = f_{big} \circ i_Y$ découle de l'égalité
$i_f^{-1} = i_T^{-1} \circ f_{big}^{-1}$, qui est claire d'après
les descriptions de ces foncteurs ci-dessus.
Nous obtenons ainsi (1).

\medskip\noindent
Le foncteur $u : X_{spaces, \etale} \to Y_{spaces, \etale}$,
$u(U \to X) = (U \times_X Y \to Y)$, comme on l'a montré, induit
un morphisme de sites et donne naissance au morphisme correspondant entre les
petits topos \'etales ; voir
Propriétés des espaces, lemme
\ref{spaces-properties-lemma-functoriality-etale-site}. La description
de l'image directe est claire.

\medskip\noindent
La partie (3) découle du fait que $\pi_X$ et $\pi_Y$ sont donnés par les
foncteurs d'inclusion, tandis que $f_{spaces, \etale}$ et $f_{big}$ sont donnés par les
foncteurs de changement de base $U \mapsto U \times_X Y$.

\medskip\noindent
L'énoncé (4) découle de (3) par précomposition avec $i_Y$.
\end{proof}

\noindent
Dans la situation du lemme, en utilisant la terminologie de
la définition \ref{definition-restriction-small-etale},
nous avons, pour $\mathcal{F}$ un faisceau sur le grand site \'etale de $Y$,
$$
(f_{big, *}\mathcal{F})|_{X_\etale} =
f_{small, *}(\mathcal{F}|_{Y_\etale}),
$$
Cette égalité découle de la commutativité du diagramme de
sites du lemme, puisque la restriction au petit site \'etale de
$Y$, resp.\ $X$, est donnée par $\pi_{Y, *}$, resp.\ $\pi_{X, *}$.
Une formule analogue faisant intervenir les images inverses et les restrictions est fausse.

\begin{lemma}
\label{lemma-composition-etale}
Soit $S$ un schéma. Étant donnés des morphismes $f : X \to Y$, $g : Y \to Z$
dans $(\textit{Spaces}/S)_\etale$, nous avons
$g_{big} \circ f_{big} = (g \circ f)_{big}$ et
$g_{small} \circ f_{small} = (g \circ f)_{small}$.
\end{lemma}

\begin{proof}
Cela découle de la description simple des foncteurs image directe
et image inverse sur les grands sites donnée par le
lemme \ref{lemma-morphism-big-etale}. Pour les foncteurs
sur les petits sites, cela découle de la description des
foncteurs image directe dans le lemme \ref{lemma-morphism-big-small-etale}.
\end{proof}

\begin{lemma}
\label{lemma-morphism-big-small-cartesian-diagram-etale}
Soit $S$ un schéma. Considérons un diagramme cartésien
$$
\xymatrix{
Y' \ar[r]_{g'} \ar[d]_{f'} & Y \ar[d]^f \\
X' \ar[r]^g & X
}
$$
dans $(\textit{Spaces}/S)_\etale$.
$i_g^{-1} \circ f_{big, *} = f'_{small, *} \circ (i_{g'})^{-1}$
et $g_{big}^{-1} \circ f_{big, *} = f'_{big, *} \circ (g'_{big})^{-1}$.
\end{lemma}

\begin{proof}
Puisque le diagramme est cartésien, nous avons, pour $U'/X'$,
$U' \times_{X'} Y' = U' \times_X Y$. Ainsi, les deux foncteurs
$i_g^{-1} \circ f_{big, *}$ et $f'_{small, *} \circ (i_{g'})^{-1}$
associent à tout faisceau $\mathcal{F}$ sur $(\textit{Spaces}/Y)_\etale$ le faisceau
$U' \mapsto \mathcal{F}(U' \times_{X'} Y')$ sur $X'_\etale$
(voir les lemmes \ref{lemma-put-in-T-etale} et \ref{lemma-morphism-big-etale}).
La seconde égalité se démontre de la même manière ou peut être
déduite du résultat très général de
Sites, lemme \ref{sites-lemma-localize-morphism}.
\end{proof}

\begin{remark}
\label{remark-change-topologies-ringed}
Les sites $(\textit{Spaces}/X)_\etale$ et $X_{spaces, \etale}$
sont munis de faisceaux structuraux. Pour le petit site \'etale, nous
l'avons vu dans Propriétés des espaces, section
\ref{spaces-properties-section-structure-sheaf}.
Le faisceau structural $\mathcal{O}$ sur le grand site \'etale
$(\textit{Spaces}/X)_\etale$ est défini en associant à un objet
$U$ les sections globales du faisceau structural de $U$.
Cela a un sens puisque $U$ est lui-même un espace algébrique
et possède donc un faisceau structural. Comme $\mathcal{O}_U$
est un faisceau sur le site \'etale de $U$, le préfaisceau $\mathcal{O}$
ainsi défini satisfait à la condition de faisceau pour les recouvrements de $U$, c'est-à-dire que
$\mathcal{O}$ est un faisceau.
Nous pouvons compléter les morphismes $i_f$, $\pi_X$, $i_X$, $f_{small}$ et
$f_{big}$ définis ci-dessus pour en faire des morphismes de sites annelés, respectivement de topos annelés.
Examinons-les successivement.
\begin{enumerate}
\item Dans le lemme \ref{lemma-put-in-T-etale}, notons $\mathcal{O}$
le faisceau structural sur $(\textit{Spaces}/X)_\etale$.
Nous avons $(i_f^{-1}\mathcal{O})(U/Y) = \mathcal{O}_U(U) = \mathcal{O}_Y(U)$
par construction.
On obtient ainsi un isomorphisme $i_f^\sharp : i_f^{-1}\mathcal{O} \to \mathcal{O}_Y$.
\item Dans le lemme \ref{lemma-at-the-bottom-etale}, il a été noté
que $i_X$ est un cas particulier de $i_f$ avec $f = \text{id}_X$ ;
nous sommes donc ramenés au cas (1).
\item Dans le lemme \ref{lemma-at-the-bottom-etale}, le morphisme
$\pi_X$ vérifie $(\pi_{X, *}\mathcal{O})(U) = \mathcal{O}(U) =
\mathcal{O}_X(U)$. Par conséquent, nous pouvons utiliser ceci pour définir
$\pi_X^\sharp : \mathcal{O}_X \to \pi_{X, *}\mathcal{O}$.
\item Dans le lemme \ref{lemma-morphism-big-small-etale},
l'extension de $f_{small}$ en un morphisme de topos annelés
a été discutée dans Propriétés des espaces, lemme
\ref{spaces-properties-lemma-morphism-ringed-topoi}.
\item Dans le lemme \ref{lemma-morphism-big-small-etale},
le foncteur $f_{big}^{-1}$ est simplement la restriction
via le foncteur d'inclusion
$(\textit{Spaces}/Y)_\etale \to (\textit{Spaces}/X)_\etale$.
Soit $\mathcal{O}_1$ le faisceau structural sur $(\textit{Spaces}/X)_\etale$
et soit $\mathcal{O}_2$ le faisceau structural sur $(\textit{Spaces}/Y)_\etale$.
Nous obtenons un isomorphisme canonique
$f_{big}^\sharp : f_{big}^{-1}\mathcal{O}_1 \to \mathcal{O}_2$.
\end{enumerate}
De plus, ces définitions sont également compatibles avec les compositions.
Nous omettons l'énoncé et la démonstration détaillés.
\end{remark}









\section{Topologie lisse}
\label{section-smooth}

\noindent
Dans cette section, nous étudions la notion de recouvrement lisse des
espaces algébriques et définissons le grand site lisse d'un
espace algébrique. Comparer avec
Topologies, section \ref{topologies-section-smooth}.

\begin{definition}
\label{definition-smooth-covering}
Soit $S$ un schéma et soit $X$ un espace algébrique sur $S$.
Un {\it recouvrement lisse de $X$} est une famille de morphismes
$\{f_i : X_i \to X\}_{i \in I}$ d'espaces algébriques sur $S$
telle que chaque $f_i$ soit lisse
et telle que
$$
|X| = \bigcup\nolimits_{i \in I} |f_i|(|X_i|),
$$
c'est-à-dire que la famille des morphismes est surjective.
\end{definition}

\noindent
C'est exactement la même notion que celle de
Topologies, définition \ref{topologies-definition-smooth-covering}.
En particulier, si $X$ et tous les $X_i$ sont des schémas, on retrouve la
notion usuelle de recouvrement lisse de schémas.

\begin{lemma}
\label{lemma-zariski-etale-smooth}
Tout recouvrement \'etale est un recouvrement lisse, et a fortiori,
tout recouvrement de Zariski est un recouvrement lisse.
\end{lemma}

\begin{proof}
Cela découle immédiatement des définitions et du fait qu'un
morphisme \'etale est lisse
(Morphismes des espaces, lemme \ref{spaces-morphisms-lemma-etale-smooth}), et
du lemme \ref{lemma-zariski-etale}.
\end{proof}

\begin{lemma}
\label{lemma-smooth}
Soit $S$ un schéma.
Soit $X$ un espace algébrique sur $S$.
\begin{enumerate}
\item Si $X' \to X$ est un isomorphisme, alors $\{X' \to X\}$
est un recouvrement lisse de $X$.
\item Si $\{X_i \to X\}_{i\in I}$ est un recouvrement lisse et si, pour chaque
$i$, nous avons un recouvrement lisse $\{X_{ij} \to X_i\}_{j\in J_i}$, alors
$\{X_{ij} \to X\}_{i \in I, j\in J_i}$ est un recouvrement lisse.
\item Si $\{X_i \to X\}_{i\in I}$ est un recouvrement lisse
et si $X' \to X$ est un morphisme d'espaces algébriques, alors
$\{X' \times_X X_i \to X'\}_{i\in I}$ est un recouvrement lisse.
\end{enumerate}
\end{lemma}

\begin{proof}
Démonstration omise.
\end{proof}

\noindent
À suivre...









\section{Topologie syntomique}
\label{section-syntomic}

\noindent
Dans cette section, nous étudions la notion de recouvrement syntomique des
espaces algébriques et définissons le grand site syntomique d'un
espace algébrique. Comparer avec
Topologies, section \ref{topologies-section-syntomic}.

\begin{definition}
\label{definition-syntomic-covering}
Soit $S$ un schéma et soit $X$ un espace algébrique sur $S$.
Un {\it recouvrement syntomique de $X$} est une famille de morphismes
$\{f_i : X_i \to X\}_{i \in I}$ d'espaces algébriques sur $S$
telle que chaque $f_i$ soit syntomique
et telle que
$$
|X| = \bigcup\nolimits_{i \in I} |f_i|(|X_i|),
$$
c'est-à-dire que la famille des morphismes est surjective.
\end{definition}

\noindent
C'est exactement la même notion que celle de
Topologies, définition \ref{topologies-definition-syntomic-covering}.
En particulier, si $X$ et tous les $X_i$ sont des schémas, on retrouve la
notion usuelle de recouvrement syntomique de schémas.

\begin{lemma}
\label{lemma-zariski-etale-smooth-syntomic}
Tout recouvrement lisse est un recouvrement syntomique, et a fortiori,
tout recouvrement \'etale ou de Zariski est un recouvrement syntomique.
\end{lemma}

\begin{proof}
Cela découle immédiatement des définitions et du fait qu'un morphisme lisse
est syntomique
(Morphismes des espaces, lemme \ref{spaces-morphisms-lemma-smooth-syntomic}),
ainsi que du lemme \ref{lemma-zariski-etale-smooth}.
\end{proof}

\begin{lemma}
\label{lemma-syntomic}
Soit $S$ un schéma.
Soit $X$ un espace algébrique sur $S$.
\begin{enumerate}
\item Si $X' \to X$ est un isomorphisme, alors $\{X' \to X\}$
est un recouvrement syntomique de $X$.
\item Si $\{X_i \to X\}_{i\in I}$ est un recouvrement syntomique et si, pour chaque
$i$, nous avons un recouvrement syntomique $\{X_{ij} \to X_i\}_{j\in J_i}$, alors
$\{X_{ij} \to X\}_{i \in I, j\in J_i}$ est un recouvrement syntomique.
\item Si $\{X_i \to X\}_{i\in I}$ est un recouvrement syntomique
et si $X' \to X$ est un morphisme d'espaces algébriques, alors
$\{X' \times_X X_i \to X'\}_{i\in I}$ est un recouvrement syntomique.
\end{enumerate}
\end{lemma}

\begin{proof}
Démonstration omise.
\end{proof}

\noindent
À suivre...










\section{Topologie fppf}
\label{section-fppf}

\noindent
Dans cette section, nous étudions la notion de recouvrement fppf d'espaces algébriques,
et définissons le grand site fppf d'un espace algébrique. Comparer avec
Topologies, section \ref{topologies-section-fppf}.

\begin{definition}
\label{definition-fppf-covering}
Soit $S$ un schéma et soit $X$ un espace algébrique sur $S$.
Un {\it recouvrement fppf de $X$} est une famille de morphismes
$\{f_i : X_i \to X\}_{i \in I}$ d'espaces algébriques sur $S$
telle que chaque $f_i$ soit plat et localement de présentation finie
et telle que
$$
|X| = \bigcup\nolimits_{i \in I} |f_i|(|X_i|),
$$
c'est-à-dire que la famille des morphismes est surjective.
\end{definition}

\noindent
C'est exactement la même notion que celle de
Topologies, définition \ref{topologies-definition-fppf-covering}.
En particulier, si $X$ et tous les $X_i$ sont des schémas, on retrouve la notion usuelle
de recouvrement fppf de schémas.

\begin{lemma}
\label{lemma-zariski-etale-smooth-syntomic-fppf}
Tout recouvrement syntomique est un recouvrement fppf, et a fortiori,
tout recouvrement lisse, \'etale ou de Zariski est un recouvrement fppf.
\end{lemma}

\begin{proof}
Cela découle des définitions et du fait qu'un morphisme syntomique
est plat et localement de présentation finie
(Morphismes des espaces, lemmes
\ref{spaces-morphisms-lemma-syntomic-locally-finite-presentation} et
\ref{spaces-morphisms-lemma-syntomic-flat}) et
du lemme \ref{lemma-zariski-etale-smooth-syntomic}.
\end{proof}

\begin{lemma}
\label{lemma-fppf}
Soit $S$ un schéma.
Soit $X$ un espace algébrique sur $S$.
\begin{enumerate}
\item Si $X' \to X$ est un isomorphisme, alors $\{X' \to X\}$
est un recouvrement fppf de $X$.
\item Si $\{X_i \to X\}_{i\in I}$ est un recouvrement fppf et si, pour chaque
$i$, nous avons un recouvrement fppf $\{X_{ij} \to X_i\}_{j\in J_i}$, alors
$\{X_{ij} \to X\}_{i \in I, j\in J_i}$ est un recouvrement fppf.
\item Si $\{X_i \to X\}_{i\in I}$ est un recouvrement fppf
et si $X' \to X$ est un morphisme d'espaces algébriques, alors
$\{X' \times_X X_i \to X'\}_{i\in I}$ est un recouvrement fppf.
\end{enumerate}
\end{lemma}

\begin{proof}
Démonstration omise.
\end{proof}

\begin{lemma}
\label{lemma-refine-fppf-schemes}
Soit $S$ un schéma et soit $X$ un espace algébrique sur $S$.
Supposons que $\mathcal{U} = \{f_i : X_i \to X\}_{i \in I}$ soit un
recouvrement fppf de $X$. Alors il existe un raffinement
$\mathcal{V} = \{g_i : T_i \to X\}$ de $\mathcal{U}$ qui est un
recouvrement fppf tel que chaque $T_i$ soit un schéma.
\end{lemma}

\begin{proof}
Démonstration omise. Indication : pour chaque $i$, choisir un schéma $T_i$ et un morphisme \'etale surjectif
$T_i \to X_i$. Vérifier ensuite que $\{T_i \to X\}$ est un recouvrement fppf.
\end{proof}

\begin{lemma}
\label{lemma-fppf-covering-surjective}
Soit $S$ un schéma.
Soit $\{f_i : X_i \to X\}_{i \in I}$ un recouvrement fppf d'espaces algébriques
sur $S$. Alors le morphisme de faisceaux
$$
\coprod X_i \longrightarrow X
$$
est surjectif.
\end{lemma}

\begin{proof}
Cela découle du
lemme \ref{spaces-lemma-surjective-flat-locally-finite-presentation} d'Espaces.
Voir également
la remarque \ref{spaces-remark-warning} d'Espaces
si la signification de ce lemme vous semble obscure.
\end{proof}

\begin{definition}
\label{definition-big-fppf-site}
Soit $S$ un schéma. Un grand site fppf {\it $(\textit{Spaces}/S)_{fppf}$}
est un site construit comme suit :
\begin{enumerate}
\item Choisissons un grand site fppf $(\Sch/S)_{fppf}$ comme dans
Topologies, section \ref{topologies-section-fppf}.
\item Prenons pour catégorie sous-jacente la catégorie $\textit{Spaces}/S$
des espaces algébriques sur $S$ (voir la discussion de la
section \ref{section-procedure} expliquant qu'il s'agit d'un ensemble).
\item Choisissons un ensemble quelconque de recouvrements comme dans
Ensembles, lemme \ref{sets-lemma-coverings-site}, en partant de la
catégorie $\textit{Spaces}/S$ et de la classe des recouvrements fppf
de la définition \ref{definition-fppf-covering}.
\end{enumerate}
\end{definition}

\noindent
Après cette définition, on peut localiser pour obtenir le site fppf
d'un espace algébrique.

\begin{definition}
\label{definition-big-small-fppf}
Soit $S$ un schéma. Soit $(\textit{Spaces}/S)_{fppf}$ comme dans
la définition \ref{definition-big-fppf-site}.
Soit $X$ un espace algébrique sur $S$, c'est-à-dire un objet de
$(\textit{Spaces}/S)_{fppf}$. Alors le grand site fppf
{\it $(\textit{Spaces}/X)_{fppf}$} de $X$
est la localisation du site $(\textit{Spaces}/S)_{fppf}$
en $X$, introduite dans Sites, section \ref{sites-section-localize}.
\end{definition}

\noindent
Nous établissons ensuite quelques relations entre les topos
associés à ces sites.

\begin{lemma}
\label{lemma-morphism-big-fppf}
Soit $S$ un schéma.
Soit $f : Y \to X$ un morphisme d'espaces algébriques sur $S$.
Le foncteur
$$
u : (\textit{Spaces}/Y)_{fppf} \longrightarrow (\textit{Spaces}/X)_{fppf},
\quad
V/Y \longmapsto V/X
$$
est cocontinu et possède un adjoint à droite continu
$$
v : (\textit{Spaces}/X)_{fppf} \longrightarrow (\textit{Spaces}/Y)_{fppf},
\quad
(U \to Y) \longmapsto (U \times_X Y \to Y).
$$
Ils induisent le même morphisme de topos
$$
f_{big} :
\Sh((\textit{Spaces}/Y)_{fppf})
\longrightarrow
\Sh((\textit{Spaces}/X)_{fppf})
$$
Nous avons $f_{big}^{-1}(\mathcal{G})(U/Y) = \mathcal{G}(U/X)$.
Nous avons $f_{big, *}(\mathcal{F})(U/X) = \mathcal{F}(U \times_X Y/Y)$.
De plus, $f_{big}^{-1}$ possède un adjoint à gauche $f_{big!}$ qui commute aux
produits fibrés et aux égaliseurs.
\end{lemma}

\begin{proof}
Le foncteur $u$ est cocontinu, continu et commute aux produits fibrés
et aux égaliseurs. Par conséquent,
les lemmes de Sites \ref{sites-lemma-when-shriek} et
\ref{sites-lemma-preserve-equalizers}
s'appliquent et nous obtenons la formule
pour $f_{big}^{-1}$ ainsi que l'existence de $f_{big!}$. De plus,
le foncteur $v$ est un adjoint à droite car, pour $U/T$ et $V/X$,
nous avons $\Mor_X(u(U), V) = \Mor_Y(U, V \times_X Y)$
comme souhaité. Nous pouvons donc appliquer
les lemmes de Sites \ref{sites-lemma-have-functor-other-way} et
\ref{sites-lemma-have-functor-other-way-morphism} pour obtenir la
formule de $f_{big, *}$.
\end{proof}

\begin{lemma}
\label{lemma-composition-fppf}
Soit $S$ un schéma. Étant donnés des morphismes $f : X \to Y$, $g : Y \to Z$
d'espaces algébriques sur $S$, nous avons
$g_{big} \circ f_{big} = (g \circ f)_{big}$.
\end{lemma}

\begin{proof}
Cela découle de la description simple des foncteurs image directe
et image inverse sur les grands sites donnée par le
lemme \ref{lemma-morphism-big-fppf}.
\end{proof}








\section{La topologie ph}
\label{section-ph}

\noindent
Dans cette section, nous définissons la topologie ph. Il s'agit de la topologie
engendrée par les recouvrements \'etales et les morphismes propres surjectifs ; voir
le lemme \ref{lemma-characterize-sheaf}.

\begin{definition}
\label{definition-ph-covering}
Soit $S$ un schéma et soit $X$ un espace algébrique sur $S$.
Un {\it recouvrement ph de $X$} est une famille
de morphismes $\{X_i \to X\}_{i \in I}$ d'espaces algébriques sur $S$
telle que chaque $f_i$ soit localement de type fini et telle que, pour tout
$U \to X$ avec $U$ affine, il existe un recouvrement ph standard
$\{U_j \to U\}_{j = 1, \ldots, m}$ qui raffine la famille
$\{X_i \times_X U \to U\}_{i \in I}$.
\end{definition}

\noindent
Autrement dit, il existe des indices $i_1, \ldots, i_m \in I$ et des
morphismes $h_j : U_j \to X_{i_j}$ tels que
$f_{i_j} \circ h_j = h \circ g_j$. Notons que, si $X$ et tous les $X_i$ sont
représentables, il s'agit de la même notion qu'un recouvrement ph de schémas défini dans
Topologies, définition \ref{topologies-definition-ph-covering}.

\begin{lemma}
\label{lemma-zariski-etale-smooth-syntomic-fppf-ph}
Tout recouvrement fppf est un recouvrement ph, et a fortiori,
tout recouvrement syntomique, lisse, \'etale ou de Zariski est un recouvrement ph.
\end{lemma}

\begin{proof}
Montrons qu'un recouvrement fppf est un recouvrement ph ; le
reste découle alors du lemme \ref{lemma-zariski-etale-smooth-syntomic-fppf}.
Soit $\{X_i \to X\}_{i \in I}$ un recouvrement fppf d'espaces algébriques
sur un schéma de base $S$. Soit $U$ un schéma affine et soit
$U \to X$ un morphisme. Nous pouvons raffiner le recouvrement fppf
$\{X_i \times_U U \to U\}_{i \in I}$ par un recouvrement fppf
$\{T_i \to U\}_{i \in I}$ où chaque $T_i$ est un schéma
(lemme \ref{lemma-refine-fppf-schemes}).
Nous pouvons alors trouver un recouvrement ph standard $\{U_j \to U\}_{j = 1, \ldots, m}$
raffinant $\{T_i \to U\}_{i \in I}$, par le
lemme \ref{more-morphisms-lemma-fppf-ph} de Compléments sur les morphismes
(et par la définition des recouvrements ph pour les schémas).
Ainsi, $\{X_i \to X\}_{i \in I}$ est un recouvrement ph par définition.
\end{proof}

\begin{lemma}
\label{lemma-surjective-proper-ph}
Soit $S$ un schéma. Soit $f : Y \to X$ un morphisme propre surjectif
d'espaces algébriques sur $S$. Alors $\{Y \to X\}$ est un recouvrement ph.
\end{lemma}

\begin{proof}
Soit $U \to X$ un morphisme, où $U$ est affine.
Par le lemme de Chow (sous la forme faible donnée dans
Cohomologie des espaces, lemme \ref{spaces-cohomology-lemma-weak-chow}),
il existe un morphisme propre surjectif de schémas
$V \to U$ qui se factorise par $Y \times_X U \to U$.
En prenant un recouvrement ouvert affine fini quelconque de $V$, nous obtenons un
recouvrement ph standard de $U$ qui raffine $\{X \times_Y U \to U\}$,
comme souhaité.
\end{proof}

\begin{lemma}
\label{lemma-ph}
Soit $S$ un schéma. Soit $X$ un espace algébrique sur $S$.
\begin{enumerate}
\item Si $X' \to X$ est un isomorphisme, alors $\{X' \to X\}$
est un recouvrement ph de $X$.
\item Si $\{X_i \to X\}_{i\in I}$ est un recouvrement ph et si, pour tout
$i$, on a un recouvrement ph $\{X_{ij} \to X_i\}_{j\in J_i}$, alors
$\{X_{ij} \to X\}_{i \in I, j\in J_i}$ est un recouvrement ph.
\item Si $\{X_i \to X\}_{i\in I}$ est un recouvrement ph
et si $X' \to X$ est un morphisme d'espaces algébriques, alors
$\{X' \times_X X_i \to X'\}_{i\in I}$ est un recouvrement ph.
\end{enumerate}
\end{lemma}

\begin{proof}
La partie (1) est claire. Considérons $g : X' \to X$ et
un recouvrement ph $\{X_i \to X\}_{i\in I}$ comme dans (3). Par le
lemme \ref{spaces-morphisms-lemma-base-change-finite-type} de Morphismes des espaces,
les morphismes $X' \times_X X_i \to X'$ sont localement de type fini.
Si $h' : Z \to X'$ est un morphisme d'un schéma affine
vers $X'$, posons $h = g \circ h' : Z \to X$. L'hypothèse
sur $\{X_i \to X\}_{i\in I}$ signifie qu'il existe un recouvrement ph standard
$\{Z_j \to Z\}_{j = 1, \ldots, n}$ et des morphismes $Z_j \to X_{i(j)}$ au-dessus de
$h$, pour certains $i(j) \in I$. Par la propriété universelle du produit fibré,
nous obtenons aussi des morphismes $Z_j \to X' \times_X X_{i(j)}$ au-dessus de $h'$.
Ainsi, $\{X' \times_X X_i \to X'\}_{i\in I}$ est un recouvrement ph.
Ceci démontre (3).

\medskip\noindent
Soient $\{X_i \to X\}_{i\in I}$ et $\{X_{ij} \to X_i\}_{j\in J_i}$ comme
dans (2). Soit $h : Z \to X$ un morphisme d'un schéma affine vers $X$.
Par hypothèse, il existe un recouvrement ph standard
$\{Z_j \to Z\}_{j = 1, \ldots, n}$ et des morphismes $h_j : Z_j \to X_{i(j)}$
au-dessus de $h$, pour certains indices $i(j) \in I$. Par hypothèse, il existe des
recouvrements ph standard
$\{Z_{j, l} \to Z_j\}_{l = 1, \ldots, n(j)}$
et des morphismes $Z_{j, l} \to X_{i(j)j(l)}$ au-dessus de
$h_j$, pour certains indices $j(l) \in J_{i(j)}$. Par le
lemme \ref{topologies-lemma-refine-by-standard-ph} de Topologies,
la famille $\{Z_{j, l} \to Z\}$ peut être raffinée par un recouvrement ph standard.
Nous en concluons que $\{X_{ij} \to X\}_{i \in I, j\in J_i}$
est un recouvrement ph.
\end{proof}

\begin{definition}
\label{definition-big-ph-site}
Soit $S$ un schéma. On appelle grand site ph {\it $(\textit{Spaces}/S)_{ph}$}
tout site construit comme suit :
\begin{enumerate}
\item Choisissons un grand site ph $(\Sch/S)_{ph}$ comme dans
Topologies, section \ref{topologies-section-ph}.
\item Prenons pour catégorie sous-jacente la catégorie $\textit{Spaces}/S$
des espaces algébriques sur $S$ (voir la discussion de la
section \ref{section-procedure}, qui explique pourquoi il s'agit d'un ensemble).
\item Choisissons un ensemble quelconque de recouvrements comme dans le
lemme \ref{sets-lemma-coverings-site} d'Ensembles, en partant de la
catégorie $\textit{Spaces}/S$ et de la classe des recouvrements ph
de la définition \ref{definition-ph-covering}.
\end{enumerate}
\end{definition}

\noindent
Après cette définition, on peut localiser pour obtenir le site ph
d'un espace algébrique.

\begin{definition}
\label{definition-big-small-ph}
Soit $S$ un schéma. Soit $(\textit{Spaces}/S)_{ph}$ comme dans la
définition \ref{definition-big-ph-site}.
Soit $X$ un espace algébrique sur $S$, c'est-à-dire un objet de
$(\textit{Spaces}/S)_{ph}$. Alors le grand site ph
{\it $(\textit{Spaces}/X)_{ph}$} de $X$
est la localisation du site $(\textit{Spaces}/S)_{ph}$
en $X$, introduite dans Sites, section \ref{sites-section-localize}.
\end{definition}

\noindent
Voici la caractérisation annoncée des faisceaux ph.

\begin{lemma}
\label{lemma-characterize-sheaf}
Soit $S$ un schéma. Soit $X$ un espace algébrique sur $S$.
Soit $\mathcal{F}$ un préfaisceau sur $(\textit{Spaces}/X)_{ph}$.
Alors $\mathcal{F}$ est un faisceau si et seulement si
\begin{enumerate}
\item $\mathcal{F}$ satisfait à la condition de faisceau pour les recouvrements \'etales, et
\item si $f : V \to U$ est un morphisme propre surjectif de
$(\textit{Spaces}/X)_{ph}$, alors
$\mathcal{F}(U)$ s'envoie bijectivement sur l'égaliseur
des deux applications $\mathcal{F}(V) \to \mathcal{F}(V \times_U V)$.
\end{enumerate}
\end{lemma}

\begin{proof}
Montrons que, si (1) et (2) sont satisfaites, alors $\mathcal{F}$ est un faisceau.
Soit $\{T_i \to T\}$ un recouvrement ph, c'est-à-dire un recouvrement dans
$(\textit{Spaces}/X)_{ph}$.
Vérifions la condition de faisceau pour ce recouvrement.
Soient $s_i \in \mathcal{F}(T_i)$ des sections dont les restrictions sont égales
sur $T_i \times_T T_{i'}$. Montrons qu'il existe une
unique section $s \in \mathcal{F}$ dont la restriction est $s_i$ sur $T_i$.
Soit $\{U_j \to T\}$ un recouvrement \'etale avec les $U_j$ affines.
D'après la propriété (1), il suffit de produire des sections $s_j \in \mathcal{F}(U_j)$
qui coïncident sur $U_j \cap U_{j'}$ afin de construire $s$.
Considérons les recouvrements ph $\{T_i \times_T U_j \to U_j\}$.
Alors les $s_{ji} = s_i|_{T_i \times_T U_j}$ sont des sections qui coïncident
sur $(T_i \times_T U_j) \times_{U_j} (T_{i'} \times_T U_j)$.
Choisissons un morphisme propre surjectif $V_j \to U_j$ et un recouvrement
ouvert affine fini $V_j = \bigcup V_{jk}$ tels que le recouvrement ph standard
$\{V_{jk} \to U_j\}$ raffine $\{T_i \times_T U_j \to U_j\}$.
Si $s_{jk} \in \mathcal{F}(V_{jk})$
désigne l'image inverse de $s_{ji}$ sur $V_{jk}$ par les
morphismes implicites, alors les $s_{jk}$ se recollent en une section
$s'_j \in \mathcal{F}(V_j)$. En utilisant encore une fois leur coïncidence sur les intersections,
nous voyons que $s'_j$ appartient à l'égaliseur des deux
applications $\mathcal{F}(V_j) \to \mathcal{F}(V_j \times_{U_j} V_j)$.
Par (2), la section $s'_j$ provient donc d'une unique section
$s_j \in \mathcal{F}(U_j)$. Nous omettons de vérifier que ces
sections $s_j$ possèdent toutes les propriétés voulues.
\end{proof}

\noindent
Nous établissons ensuite quelques relations entre les topos
associés à ces sites.

\begin{lemma}
\label{lemma-morphism-big-ph}
Soit $S$ un schéma.
Soit $f : Y \to X$ un morphisme d'espaces algébriques sur $S$.
Le foncteur
$$
u : (\textit{Spaces}/Y)_{ph} \longrightarrow (\textit{Spaces}/X)_{ph},
\quad
V/Y \longmapsto V/X
$$
est cocontinu et possède un adjoint à droite continu
$$
v : (\textit{Spaces}/X)_{ph} \longrightarrow (\textit{Spaces}/Y)_{ph},
\quad
(U \to Y) \longmapsto (U \times_X Y \to Y).
$$
Ils induisent le même morphisme de topos
$$
f_{big} :
\Sh((\textit{Spaces}/Y)_{ph})
\longrightarrow
\Sh((\textit{Spaces}/X)_{ph})
$$
Nous avons $f_{big}^{-1}(\mathcal{G})(U/Y) = \mathcal{G}(U/X)$.
Nous avons $f_{big, *}(\mathcal{F})(U/X) = \mathcal{F}(U \times_X Y/Y)$.
De plus, $f_{big}^{-1}$ possède un adjoint à gauche $f_{big!}$ qui commute aux
produits fibrés et aux égaliseurs.
\end{lemma}

\begin{proof}
Le foncteur $u$ est cocontinu, continu et commute aux produits fibrés
et aux égaliseurs. Par conséquent,
les lemmes de Sites \ref{sites-lemma-when-shriek} et
\ref{sites-lemma-preserve-equalizers}
s'appliquent et nous en déduisons la formule
pour $f_{big}^{-1}$ ainsi que l'existence de $f_{big!}$. De plus,
le foncteur $v$ est un adjoint à droite car, pour $U/T$ et $V/X$,
nous avons $\Mor_X(u(U), V) = \Mor_Y(U, V \times_X Y)$,
comme souhaité. Nous pouvons donc appliquer
les lemmes de Sites \ref{sites-lemma-have-functor-other-way} et
\ref{sites-lemma-have-functor-other-way-morphism} pour obtenir la
formule de $f_{big, *}$.
\end{proof}

\begin{lemma}
\label{lemma-composition-ph}
Soit $S$ un schéma. Étant donnés des morphismes $f : X \to Y$, $g : Y \to Z$
d'espaces algébriques sur $S$, nous avons
$g_{big} \circ f_{big} = (g \circ f)_{big}$.
\end{lemma}

\begin{proof}
Cela découle de la description simple des foncteurs image directe
et image inverse sur les grands sites donnée par le
lemme \ref{lemma-morphism-big-ph}.
\end{proof}

\begin{lemma}
\label{lemma-cech-enough}
Soit $S$ un schéma. Soit $X$ un espace algébrique sur $S$.
Soit $P$ une propriété des objets de $(\textit{Spaces}/X)_{fppf}$
telle que, pour tout recouvrement $\{U_i \to U\}$ de
$(\textit{Spaces}/X)_{fppf}$, l'implication suivante soit vérifiée :
$$
P(U_{i_0} \times_U \ldots \times_U U_{i_p})
\text{ pour tous }
p \geq 0,\ i_0, \ldots, i_p \in I
\Rightarrow P(U)
$$
Si $P(U)$ pour tout $U$ affine, plat et localement de présentation finie sur $X$,
alors $P(X)$.
\end{lemma}

\begin{proof}
Soit $U$ un espace algébrique séparé et localement de présentation finie sur $X$.
Nous pouvons choisir un recouvrement \'etale $\{U_i \to U\}_{i \in I}$ avec $V_i$
affine. Comme $U$ est séparé, nous en concluons que
$U_{i_0} \times_U \ldots \times_U U_{i_p}$ est toujours affine.
Ainsi, $P(U_{i_0} \times_U \ldots \times_U U_{i_p})$ est toujours vraie.
Par conséquent, $P(U)$ est vraie. Choisissons un schéma $U$ qui est une réunion disjointe de
schémas affines et un morphisme \'etale surjectif $U \to X$.
Alors $U \times_X \ldots \times_X U$ (avec $p + 1$ facteurs)
est un espace algébrique séparé
et \'etale sur $X$. D'après ce qui précède, $P(U \times_X \ldots \times_X U)$ est donc vraie.
Nous en concluons que $P(X)$ est vraie.
\end{proof}












\section{Topologie fpqc}
\label{section-fpqc}

\noindent
Nous étudions brièvement la notion de recouvrement fpqc d'espaces algébriques.
Comparer avec
Topologies, section \ref{topologies-section-fpqc}.
Nous montrerons dans
Descente sur les espaces,
proposition \ref{spaces-descent-proposition-fpqc-descent-quasi-coherent},
que les faisceaux quasi-cohérents descendent le long de tels recouvrements.

\begin{definition}
\label{definition-fpqc-covering}
Soit $S$ un schéma et soit $X$ un espace algébrique sur $S$.
Un {\it recouvrement fpqc de $X$} est une famille de morphismes
$\{f_i : X_i \to X\}_{i \in I}$ d'espaces algébriques
telle que chaque $f_i$ soit plat et que, pour tout schéma affine
$Z$ et tout morphisme $h : Z \to X$, il existe un recouvrement fpqc standard
$\{g_j : Z_j \to Z\}_{j = 1, \ldots, m}$ qui raffine la famille
$\{X_i \times_X Z \to Z\}_{i \in I}$.
\end{definition}

\noindent
Autrement dit, il existe des indices $i_1, \ldots, i_m \in I$ et des
morphismes $h_j : Z_j \to X_{i_j}$ tels que
$f_{i_j} \circ h_j = h \circ g_j$. Remarquons que, si $X$ et tous les $X_i$ sont
représentables, cela revient à un recouvrement fpqc de schémas d'après le
lemme \ref{topologies-lemma-fpqc-covering-affines-mapping-in} de Topologies.

\begin{lemma}
\label{lemma-zariski-etale-smooth-syntomic-fppf-fpqc}
Tout recouvrement fppf est un recouvrement fpqc et a fortiori,
tout recouvrement syntomique, lisse, \'etale ou de Zariski est un recouvrement fpqc.
\end{lemma}

\begin{proof}
Montrons qu'un recouvrement fppf est un recouvrement fpqc ; le
reste découle alors du
lemme \ref{lemma-zariski-etale-smooth-syntomic-fppf}.
Soit $\{f_i : U_i \to U\}_{i \in I}$ un
recouvrement fppf d'espaces algébriques sur $S$.
Par définition, les $f_i$ sont plats, ce qui vérifie la première
condition de la définition \ref{definition-fpqc-covering}. Pour vérifier la
seconde, soit $V \to U$ un morphisme avec $V$ affine.
Nous pouvons choisir un recouvrement \'etale $\{V_{ij} \to V \times_U U_i\}$
avec les $V_{ij}$ affines. Alors les composés
$f_{ij} : V_{ij} \to V \times_U U_i \to V$ sont plats et localement de
présentation finie, chacun étant composé de morphismes possédant ces propriétés
(Morphismes des espaces, lemmes
\ref{spaces-morphisms-lemma-composition-finite-presentation},
\ref{spaces-morphisms-lemma-composition-flat},
\ref{spaces-morphisms-lemma-etale-flat}, et
\ref{spaces-morphisms-lemma-etale-locally-finite-presentation}).
Ces morphismes sont donc ouverts
(Morphismes des espaces, lemme \ref{spaces-morphisms-lemma-fppf-open}),
et nous voyons que
$|V| = \bigcup_{i \in I} \bigcup_{j \in J_i} f_{ij}(|V_{ij}|)$
est un recouvrement ouvert de $|V|$.
Comme $|V|$ est quasi-compact, ce recouvrement possède un
raffinement fini.
Disons que $V_{i_1j_1}, \ldots, V_{i_Nj_N}$ conviennent.
Alors $\{V_{i_kj_k} \to V\}_{k = 1, \ldots, N}$ est
un recouvrement fpqc standard de $V$ qui raffine la famille
$\{U_i \times_U V \to V\}$.
Ceci achève la démonstration.
\end{proof}

\begin{lemma}
\label{lemma-fpqc}
Soit $S$ un schéma.
Soit $X$ un espace algébrique sur $S$.
\begin{enumerate}
\item Si $X' \to X$ est un isomorphisme, alors $\{X' \to X\}$
est un recouvrement fpqc de $X$.
\item Si $\{X_i \to X\}_{i\in I}$ est un recouvrement fpqc et si, pour tout
$i$, on a un recouvrement fpqc $\{X_{ij} \to X_i\}_{j\in J_i}$, alors
$\{X_{ij} \to X\}_{i \in I, j\in J_i}$ est un recouvrement fpqc.
\item Si $\{X_i \to X\}_{i\in I}$ est un recouvrement fpqc
et si $X' \to X$ est un morphisme d'espaces algébriques, alors
$\{X' \times_X X_i \to X'\}_{i\in I}$ est un recouvrement fpqc.
\end{enumerate}
\end{lemma}

\begin{proof}
La partie (1) est claire. Considérons $g : X' \to X$ et
un recouvrement fpqc $\{X_i \to X\}_{i\in I}$ comme dans (3). Par le
lemme \ref{spaces-morphisms-lemma-base-change-flat} de Morphismes des espaces,
les morphismes $X' \times_X X_i \to X'$
sont plats. Si $h' : Z \to X'$ est un morphisme d'un schéma affine
vers $X'$, posons $h = g \circ h' : Z \to X$. L'hypothèse
sur $\{X_i \to X\}_{i\in I}$ signifie qu'il existe un recouvrement fpqc standard
$\{Z_j \to Z\}_{j = 1, \ldots, n}$ et des morphismes $Z_j \to X_{i(j)}$ au-dessus de
$h$, pour certains $i(j) \in I$. Par la propriété universelle du produit fibré,
nous obtenons aussi des morphismes $Z_j \to X' \times_X X_{i(j)}$ au-dessus de $h'$.
Ainsi, $\{X' \times_X X_i \to X'\}_{i\in I}$ est un recouvrement fpqc.
Ceci démontre (3).

\medskip\noindent
Soient $\{X_i \to X\}_{i\in I}$ et $\{X_{ij} \to X_i\}_{j\in J_i}$ comme
dans (2). Soit $h : Z \to X$ un morphisme d'un schéma affine vers $X$.
Par hypothèse, il existe un recouvrement fpqc standard
$\{Z_j \to Z\}_{j = 1, \ldots, n}$ et des morphismes $h_j : Z_j \to X_{i(j)}$
au-dessus de $h$, pour certains indices $i(j) \in I$. Par hypothèse, il existe des
recouvrements fpqc standard
$\{Z_{j, l} \to Z_j\}_{l = 1, \ldots, n(j)}$
et des morphismes $Z_{j, l} \to X_{i(j)j(l)}$ au-dessus de
$h_j$, pour certains indices $j(l) \in J_{i(j)}$. Par le
lemme \ref{topologies-lemma-fpqc-affine-axioms} de Topologies,
la famille $\{Z_{j, l} \to Z\}$ est un recouvrement fpqc standard.
Nous en concluons que $\{X_{ij} \to X\}_{i \in I, j\in J_i}$
est un recouvrement fpqc.
\end{proof}

\begin{lemma}
\label{lemma-recognize-fpqc-covering}
Soit $S$ un schéma et soit $X$ un espace algébrique sur $S$.
Supposons que $\{f_i : X_i \to X\}_{i \in I}$ soit une famille de morphismes
d'espaces algébriques de but $X$. Soit $U \to X$ un morphisme
\'etale surjectif d'un schéma vers $X$. Alors
$\{f_i : X_i \to X\}_{i \in I}$ est un recouvrement fpqc de $X$ si et seulement
si $\{U \times_X X_i \to U\}_{i \in I}$ est un recouvrement fpqc de $U$.
\end{lemma}

\begin{proof}
Si $\{X_i \to X\}_{i \in I}$ est un recouvrement fpqc, il en va de même de
$\{U \times_X X_i \to U\}_{i \in I}$ par le lemme \ref{lemma-fpqc}.
Supposons que $\{U \times_X X_i \to U\}_{i \in I}$ soit un recouvrement fpqc.
Soit $h : Z \to X$ un morphisme d'un schéma affine vers $X$.
Alors $U \times_X Z \to Z$ est un morphisme \'etale surjectif
de schémas, donc en particulier ouvert. Nous pouvons donc trouver un nombre fini d'ouverts affines
$W_1, \ldots, W_t$ de $U \times_X Z$ dont les images recouvrent $Z$.
Pour chaque $j$, nous pouvons appliquer la condition selon laquelle
$\{U \times_X X_i \to U\}_{i \in I}$ est un recouvrement fpqc
au morphisme $W_j \to U$, et obtenir un recouvrement fpqc standard
$\{W_{jl} \to W_j\}$ qui raffine $\{W_j \times_X X_i \to W_j\}_{i \in I}$.
Ainsi, $\{W_{jl} \to Z\}$ est un recouvrement fpqc standard de $Z$
(voir le
lemme \ref{topologies-lemma-fpqc-affine-axioms} de Topologies)
qui raffine $\{Z \times_X X_i \to X\}$, ce qui achève la démonstration.
\end{proof}

\begin{lemma}
\label{lemma-refine-fpqc-schemes}
Soit $S$ un schéma et soit $X$ un espace algébrique sur $S$.
Supposons que $\mathcal{U} = \{f_i : X_i \to X\}_{i \in I}$ soit un
recouvrement fpqc de $X$. Il existe alors un raffinement
$\mathcal{V} = \{g_i : T_i \to X\}$ de $\mathcal{U}$ qui est un
recouvrement fpqc tel que chaque $T_i$ soit un schéma.
\end{lemma}

\begin{proof}
Démonstration omise. Indication : pour chaque $i$, choisir un schéma $T_i$ et un morphisme \'etale
surjectif $T_i \to X_i$. Vérifier ensuite que $\{T_i \to X\}$ est un recouvrement fpqc.
\end{proof}

\noindent
À suivre...





\begin{multicols}{2}[\section{Autres chapitres}]
\noindent
Préliminaires
\begin{enumerate}
\item \hyperref[introduction-section-phantom]{Introduction}
\item \hyperref[conventions-section-phantom]{Conventions}
\item \hyperref[sets-section-phantom]{Théorie des ensembles}
\item \hyperref[categories-section-phantom]{Catégories}
\item \hyperref[topology-section-phantom]{Topologie}
\item \hyperref[sheaves-section-phantom]{Faisceaux sur les espaces}
\item \hyperref[sites-section-phantom]{Sites et faisceaux}
\item \hyperref[stacks-section-phantom]{Champs}
\item \hyperref[fields-section-phantom]{Corps}
\item \hyperref[algebra-section-phantom]{Algèbre commutative}
\item \hyperref[brauer-section-phantom]{Groupes de Brauer}
\item \hyperref[homology-section-phantom]{Algèbre homologique}
\item \hyperref[derived-section-phantom]{Catégories dérivées}
\item \hyperref[simplicial-section-phantom]{Méthodes simpliciales}
\item \hyperref[more-algebra-section-phantom]{Compléments d'algèbre}
\item \hyperref[smoothing-section-phantom]{Lissage des morphismes d'anneaux}
\item \hyperref[modules-section-phantom]{Faisceaux de modules}
\item \hyperref[sites-modules-section-phantom]{Modules sur les sites}
\item \hyperref[injectives-section-phantom]{Objets injectifs}
\item \hyperref[cohomology-section-phantom]{Cohomologie des faisceaux}
\item \hyperref[sites-cohomology-section-phantom]{Cohomologie sur les sites}
\item \hyperref[dga-section-phantom]{Algèbre différentielle graduée}
\item \hyperref[dpa-section-phantom]{Algèbre à puissances divisées}
\item \hyperref[sdga-section-phantom]{Faisceaux différentiels gradués}
\item \hyperref[hypercovering-section-phantom]{Hyperrecouvrements}
\end{enumerate}
Schémas
\begin{enumerate}
\setcounter{enumi}{25}
\item \hyperref[schemes-section-phantom]{Schémas}
\item \hyperref[constructions-section-phantom]{Constructions de schémas}
\item \hyperref[properties-section-phantom]{Propriétés des schémas}
\item \hyperref[morphisms-section-phantom]{Morphismes de schémas}
\item \hyperref[coherent-section-phantom]{Cohomologie des schémas}
\item \hyperref[divisors-section-phantom]{Diviseurs}
\item \hyperref[limits-section-phantom]{Limites de schémas}
\item \hyperref[varieties-section-phantom]{Variétés}
\item \hyperref[topologies-section-phantom]{Topologies sur les schémas}
\item \hyperref[descent-section-phantom]{Descente}
\item \hyperref[perfect-section-phantom]{Catégories dérivées de schémas}
\item \hyperref[more-morphisms-section-phantom]{Compléments sur les morphismes}
\item \hyperref[flat-section-phantom]{Compléments sur la platitude}
\item \hyperref[groupoids-section-phantom]{Schémas en groupoïdes}
\item \hyperref[more-groupoids-section-phantom]{Compléments sur les schémas en groupoïdes}
\item \hyperref[etale-section-phantom]{Morphismes étales de schémas}
\end{enumerate}
Thèmes de théorie des schémas
\begin{enumerate}
\setcounter{enumi}{41}
\item \hyperref[chow-section-phantom]{Homologie de Chow}
\item \hyperref[intersection-section-phantom]{Théorie de l'intersection}
\item \hyperref[pic-section-phantom]{Schémas de Picard des courbes}
\item \hyperref[weil-section-phantom]{Théories cohomologiques de Weil}
\item \hyperref[adequate-section-phantom]{Modules adéquats}
\item \hyperref[dualizing-section-phantom]{Complexes dualisants}
\item \hyperref[duality-section-phantom]{Dualité pour les schémas}
\item \hyperref[discriminant-section-phantom]{Discriminants et différentes}
\item \hyperref[derham-section-phantom]{Cohomologie de de Rham}
\item \hyperref[local-cohomology-section-phantom]{Cohomologie locale}
\item \hyperref[algebraization-section-phantom]{Géométrie algébrique et formelle}
\item \hyperref[curves-section-phantom]{Courbes algébriques}
\item \hyperref[resolve-section-phantom]{Résolution des surfaces}
\item \hyperref[models-section-phantom]{Réduction semi-stable}
\item \hyperref[functors-section-phantom]{Foncteurs et morphismes}
\item \hyperref[equiv-section-phantom]{Catégories dérivées de variétés}
\item \hyperref[pione-section-phantom]{Groupes fondamentaux de schémas}
\item \hyperref[etale-cohomology-section-phantom]{Cohomologie étale}
\item \hyperref[crystalline-section-phantom]{Cohomologie cristalline}
\item \hyperref[proetale-section-phantom]{Cohomologie pro-étale}
\item \hyperref[relative-cycles-section-phantom]{Cycles relatifs}
\item \hyperref[more-etale-section-phantom]{Compléments de cohomologie étale}
\item \hyperref[trace-section-phantom]{La formule des traces}
\end{enumerate}
Espaces algébriques
\begin{enumerate}
\setcounter{enumi}{64}
\item \hyperref[spaces-section-phantom]{Espaces algébriques}
\item \hyperref[spaces-properties-section-phantom]{Propriétés des espaces algébriques}
\item \hyperref[spaces-morphisms-section-phantom]{Morphismes d'espaces algébriques}
\item \hyperref[decent-spaces-section-phantom]{Espaces algébriques décents}
\item \hyperref[spaces-cohomology-section-phantom]{Cohomologie des espaces algébriques}
\item \hyperref[spaces-limits-section-phantom]{Limites d'espaces algébriques}
\item \hyperref[spaces-divisors-section-phantom]{Diviseurs sur les espaces algébriques}
\item \hyperref[spaces-over-fields-section-phantom]{Espaces algébriques sur des corps}
\item \hyperref[spaces-topologies-section-phantom]{Topologies sur les espaces algébriques}
\item \hyperref[spaces-descent-section-phantom]{Descente et espaces algébriques}
\item \hyperref[spaces-perfect-section-phantom]{Catégories dérivées d'espaces}
\item \hyperref[spaces-more-morphisms-section-phantom]{Compléments sur les morphismes d'espaces}
\item \hyperref[spaces-flat-section-phantom]{Platitude sur les espaces algébriques}
\item \hyperref[spaces-groupoids-section-phantom]{Groupoïdes dans les espaces algébriques}
\item \hyperref[spaces-more-groupoids-section-phantom]{Compléments sur les groupoïdes dans les espaces}
\item \hyperref[bootstrap-section-phantom]{Amorçage}
\item \hyperref[spaces-pushouts-section-phantom]{Sommes amalgamées d'espaces algébriques}
\end{enumerate}
Thèmes de géométrie
\begin{enumerate}
\setcounter{enumi}{81}
\item \hyperref[spaces-chow-section-phantom]{Groupes de Chow des espaces}
\item \hyperref[groupoids-quotients-section-phantom]{Quotients de groupoïdes}
\item \hyperref[spaces-more-cohomology-section-phantom]{Compléments sur la cohomologie des espaces}
\item \hyperref[spaces-simplicial-section-phantom]{Espaces simpliciaux}
\item \hyperref[spaces-duality-section-phantom]{Dualité pour les espaces}
\item \hyperref[formal-spaces-section-phantom]{Espaces algébriques formels}
\item \hyperref[restricted-section-phantom]{Algébrisation des espaces formels}
\item \hyperref[spaces-resolve-section-phantom]{Résolution des surfaces revisitée}
\end{enumerate}
Théorie des déformations
\begin{enumerate}
\setcounter{enumi}{89}
\item \hyperref[formal-defos-section-phantom]{Théorie formelle des déformations}
\item \hyperref[defos-section-phantom]{Théorie des déformations}
\item \hyperref[cotangent-section-phantom]{Le complexe cotangent}
\item \hyperref[examples-defos-section-phantom]{Problèmes de déformation}
\end{enumerate}
Champs algébriques
\begin{enumerate}
\setcounter{enumi}{93}
\item \hyperref[algebraic-section-phantom]{Champs algébriques}
\item \hyperref[examples-stacks-section-phantom]{Exemples de champs}
\item \hyperref[stacks-sheaves-section-phantom]{Faisceaux sur les champs algébriques}
\item \hyperref[criteria-section-phantom]{Critères de représentabilité}
\item \hyperref[artin-section-phantom]{Axiomes d'Artin}
\item \hyperref[quot-section-phantom]{Espaces de Quot et de Hilbert}
\item \hyperref[stacks-properties-section-phantom]{Propriétés des champs algébriques}
\item \hyperref[stacks-morphisms-section-phantom]{Morphismes de champs algébriques}
\item \hyperref[stacks-limits-section-phantom]{Limites de champs algébriques}
\item \hyperref[stacks-cohomology-section-phantom]{Cohomologie des champs algébriques}
\item \hyperref[stacks-perfect-section-phantom]{Catégories dérivées de champs}
\item \hyperref[stacks-introduction-section-phantom]{Introduction aux champs algébriques}
\item \hyperref[stacks-more-morphisms-section-phantom]{Compléments sur les morphismes de champs}
\item \hyperref[stacks-geometry-section-phantom]{La géométrie des champs}
\end{enumerate}
Thèmes de théorie des espaces de modules
\begin{enumerate}
\setcounter{enumi}{107}
\item \hyperref[moduli-section-phantom]{Champs de modules}
\item \hyperref[moduli-curves-section-phantom]{Espaces de modules de courbes}
\end{enumerate}
Divers
\begin{enumerate}
\setcounter{enumi}{109}
\item \hyperref[examples-section-phantom]{Exemples}
\item \hyperref[exercises-section-phantom]{Exercices}
\item \hyperref[guide-section-phantom]{Guide bibliographique}
\item \hyperref[desirables-section-phantom]{Desiderata}
\item \hyperref[coding-section-phantom]{Style de programmation}
\item \hyperref[obsolete-section-phantom]{Obsolète}
\item \hyperref[fdl-section-phantom]{Licence de documentation libre GNU}
\item \hyperref[index-section-phantom]{Index généré automatiquement}
\end{enumerate}
\end{multicols}



\bibliography{my}
\bibliographystyle{amsalpha}

\end{document}
